\section{GDPval：把输出与行业专家直接比较}

\subsection{从“能完成”转向“交付是否够好”}

\videofigure{figures/economic-value.jpg}{GDPval 不只问 agent 是否完成任务，而是让专家比较 AI 与同行专业人员的真实工作产品。}{00:28:06--00:28:38}

\videofigure{figures/gdpval-sectors.jpg}{任务覆盖九个高 GDP 行业和 44 个职业，试图代表经济中常见的知识工作。}{00:28:38--00:30:02}

\videofigure{figures/gdpval-task-examples.jpg}{任务产物包括 CAD、投资分析、护理图像判断、幻灯片、表格和运营报告。}{00:30:02--00:34:16}

与代码 benchmark 不同，GDPval 的输出往往没有唯一答案。评估采用行业专家的 pairwise preference：在不知道来源的情况下比较 AI 产物与人类专家产物，并判断胜、负或平局。

\begin{importantbox}{GDPval 测的是产品质量，不是隐藏思维过程}
专家看到的是最终交付物，因此指标更接近客户体验；但它无法直接说明 agent 的过程是否可靠、是否引用了错误事实，或是否在另一输入上会崩溃。
\end{importantbox}

\subsection{数据构造与任务特征}

\videofigure{figures/gdpval-construction.jpg}{GDPval 先按 GDP 选择行业，再选职业与典型工作活动，采用自上而下的代表性构造。}{00:34:16--00:35:08}

\videofigure{figures/gdpval-task-characteristics.jpg}{单个任务的人类平均完成时间约七小时，任务多模态、依赖参考文件，且多数没有唯一正确答案。}{00:35:08--00:36:08}

这类任务比 METR 软件任务更接近办公室工作，却也更难自动评分。高质量评估需要领域专家、明确 rubric、盲测和一致性检查，成本远高于运行单元测试。

\begin{warningbox}{偏好评估会受风格与呈现影响}
排版漂亮、语言流畅的报告可能掩盖事实错误；专业但朴素的人类产物也可能被低估。应同时记录 accuracy、instruction following、formatting 和 severity，而非只看总体偏好。
\end{warningbox}

\subsection{模型进步不是同一条线}

\videofigure{figures/gdpval-improvement.jpg}{GDPval 上的 pairwise win rate 随模型代际显著提升，但呈阶梯式而非平滑线性增长。}{00:36:08--00:38:18}

\videofigure{figures/model-specific-strengths.jpg}{不同模型呈现明显专长：有的擅长美学与文档，有的更擅长计算、指令遵循或复杂工具。}{00:38:18--00:39:08}

模型排名依赖任务组合。一个模型在 slides、PDF 与视觉布局上占优，另一个可能在表格计算、文字准确性或代码工具上更强。因此企业部署应按任务路由，而非假设单一模型统治所有职业。

\begin{knowledgebox}{Model routing 是评估的直接产物}
若 benchmark 保留职业、模态和错误类型的细粒度标签，就可以训练路由器：按任务选择模型、工具链和审查强度，而不仅用于发布一个总榜。
\end{knowledgebox}

\subsection{失败严重度与经济成本}

\videofigure{figures/gdpval-failure-modes.jpg}{常见问题包括指令遵循、格式和准确性错误，不同模型的错误构成不同。}{00:39:08--00:40:08}

\videofigure{figures/failure-severity.jpg}{许多失败属于“可接受但不够好”，真正灾难性错误比例较低但不能忽略。}{00:40:08--00:41:08}

\videofigure{figures/economic-speed-cost.jpg}{考虑“尝试多次再由人修复”后，模型可能带来速度和成本收益，但收益取决于重试与审查成本。}{00:41:08--00:42:02}

可用简化期望成本模型表示：

$$
C(n)=nC_{\mathrm{AI}}+
\bigl(1-P_n(\text{acceptable})\bigr)C_{\mathrm{repair}}+C_{\mathrm{review}}.
$$

\begin{itemize}
    \item $n$：生成或重试次数；
    \item $C_{\mathrm{AI}}$：每次模型运行成本；
    \item $P_n$：从 $n$ 次尝试中得到可接受结果的概率；
    \item $C_{\mathrm{repair}}$：专家修复失败产物的成本；
    \item $C_{\mathrm{review}}$：无论成功与否都需要的审核成本。
\end{itemize}

\begin{importantbox}{经济收益来自“AI + 人”的工作流}
若审查和修复很便宜，即使单次成功率不高也可能有价值；若错误难以发现或后果严重，人工审核成本会抵消模型速度优势。
\end{importantbox}

\videofigure{figures/performance-variation.jpg}{GDPval 表现随行业、任务时长与模态变化，短任务和部分行业更接近专家水平。}{00:42:02--00:45:48}

\subsection{本章小结}

GDPval 将评估从“答对一道题”推进到“交付一个可比较的专业产物”。它揭示模型专长、错误严重度和人机协作成本，但仍受上下文和专家偏好影响。

